\section{Requisiti}
I requisiti individuati sono classificati e identificati da un codice univoco nel formato \textbf{R[ID]-[Tipo]-[Importanza]}:
\begin{itemize}
	\item \textbf{ID}: Identificativo numerico progressivo.
	\item \textbf{Tipo}: \textbf{F} (Funzionale), \textbf{Q} (Qualità), \textbf{V} (Vincolo), \textbf{P} (Prestazionale).
	\item \textbf{Importanza}: \textbf{O} (Obbligatorio), \textbf{D} (Desiderabile), \textbf{Op} (Opzionale).
\end{itemize}

La classificazione dei requisiti è stata effettuata sulla base delle
informazioni raccolte dal capitolato, dai verbali esterni con la
proponente, dall'analisi del dominio applicativo e da decisioni
interne al team di sviluppo (indicate come "Interno").

\subsection{Requisiti funzionali}
\renewcommand{\arraystretch}{1.5}
\begin{longtable}{|p{2.5cm}|p{2.2cm}|p{6.8cm}|p{2.7cm}|}
	\caption{Requisiti Funzionali} \\
	\hline
	\rowcolor{primarycolor!10}
	\textbf{Codice} & \textbf{Importanza} & \textbf{Descrizione} & \textbf{UC, Fonti} \\
	\hline
	\endfirsthead
	\hline
	\rowcolor{primarycolor!10}
	\textbf{Codice} & \textbf{Importanza} & \textbf{Descrizione} & \textbf{UC, Fonti} \\
	\hline
	\endhead

	% ---- Scrittura e clipboard ----
	\textbf{R1-F-O} & Obbligatorio & Il sistema deve permettere la scrittura libera di testo nell'editor tramite tastiera. & \hyperlink{uc1}{UC1}, Capitolato \\ \hline
	\textbf{R2-F-O} & Obbligatorio & Il sistema deve permettere la copia di una porzione di testo selezionata negli appunti di sistema. & \hyperlink{uc2}{UC2}, Analisi del dominio \\ \hline
	\textbf{R3-F-O} & Obbligatorio & Il sistema deve permettere di tagliare una porzione di testo selezionata, rimuovendola dall'editor e copiandola negli appunti. & \hyperlink{uc3}{UC3}, Analisi del dominio \\ \hline
	\textbf{R4-F-O} & Obbligatorio & Il sistema deve permettere di incollare il contenuto degli appunti nella posizione corrente del cursore. & \hyperlink{uc4}{UC4}, Analisi del dominio \\ \hline

	% ---- Formattazione inline ----
	\textbf{R5-F-O} & Obbligatorio & Il sistema deve permettere l'attivazione della formattazione grassetto per il testo digitato successivamente. & \hyperlink{uc5}{UC5}, Capitolato, VE-5.1 \\ \hline
	\textbf{R6-F-O} & Obbligatorio & Il sistema deve permettere l'applicazione della formattazione grassetto a una porzione di testo già selezionata. & \hyperlink{uc6}{UC6}, Capitolato, VE-5.1 \\ \hline
	\textbf{R7-F-O} & Obbligatorio & Il sistema deve permettere la disattivazione della formattazione grassetto durante la digitazione. & \hyperlink{uc7}{UC7}, Capitolato \\ \hline
	\textbf{R8-F-O} & Obbligatorio & Il sistema deve permettere la rimozione della formattazione grassetto da una porzione di testo selezionata. & \hyperlink{uc8}{UC8}, Capitolato \\ \hline
	\textbf{R9-F-O} & Obbligatorio & Il sistema deve permettere l'attivazione della formattazione corsivo per il testo digitato successivamente. & \hyperlink{uc9}{UC9}, Capitolato, VE-5.1 \\ \hline
	\textbf{R10-F-O} & Obbligatorio & Il sistema deve permettere l'applicazione della formattazione corsivo a una porzione di testo già selezionata. & \hyperlink{uc10}{UC10}, Capitolato, VE-5.1 \\ \hline
	\textbf{R11-F-O} & Obbligatorio & Il sistema deve permettere la disattivazione della formattazione corsivo durante la digitazione. & \hyperlink{uc11}{UC11}, Capitolato \\ \hline
	\textbf{R12-F-O} & Obbligatorio & Il sistema deve permettere la rimozione della formattazione corsivo da una porzione di testo selezionata. & \hyperlink{uc12}{UC12}, Capitolato \\ \hline
	\textbf{R13-F-O} & Obbligatorio & Il sistema deve permettere l'attivazione della formattazione sottolineato per il testo digitato successivamente. & \hyperlink{uc13}{UC13}, Capitolato, VE-5.1 \\ \hline
	\textbf{R14-F-O} & Obbligatorio & Il sistema deve permettere l'applicazione della formattazione sottolineato a una porzione di testo già selezionata. & \hyperlink{uc14}{UC14}, Capitolato, VE-5.1 \\ \hline
	\textbf{R15-F-O} & Obbligatorio & Il sistema deve permettere la disattivazione della formattazione sottolineato durante la digitazione. & \hyperlink{uc15}{UC15}, Capitolato \\ \hline
	\textbf{R16-F-O} & Obbligatorio & Il sistema deve permettere la rimozione della formattazione sottolineato da una porzione di testo selezionata. & \hyperlink{uc16}{UC16}, Capitolato \\ \hline
	\textbf{R17-F-O} & Obbligatorio & Il sistema deve permettere l'attivazione della formattazione barrato per il testo digitato successivamente. & \hyperlink{uc17}{UC17}, Capitolato, VE-5.1 \\ \hline
	\textbf{R18-F-O} & Obbligatorio & Il sistema deve permettere l'applicazione della formattazione barrato a una porzione di testo già selezionata. & \hyperlink{uc18}{UC18}, Capitolato, VE-5.1 \\ \hline
	\textbf{R19-F-O} & Obbligatorio & Il sistema deve permettere la disattivazione della formattazione barrato durante la digitazione. & \hyperlink{uc19}{UC19}, Capitolato \\ \hline
	\textbf{R20-F-O} & Obbligatorio & Il sistema deve permettere la rimozione della formattazione barrato da una porzione di testo selezionata. & \hyperlink{uc20}{UC20}, Capitolato \\ \hline
	\textbf{R21-F-O} & Obbligatorio & Il sistema deve permettere l'attivazione della formattazione citazione (blockquote) per il testo digitato successivamente. & \hyperlink{uc21}{UC21}, Capitolato, VE-5.1 \\ \hline
	\textbf{R22-F-O} & Obbligatorio & Il sistema deve permettere l'applicazione della formattazione citazione a una porzione di testo già selezionata. & \hyperlink{uc22}{UC22}, Capitolato, VE-5.1 \\ \hline
	\textbf{R23-F-O} & Obbligatorio & Il sistema deve permettere la disattivazione della formattazione citazione durante la digitazione. & \hyperlink{uc23}{UC23}, Capitolato \\ \hline
	\textbf{R24-F-O} & Obbligatorio & Il sistema deve permettere la rimozione della formattazione citazione da una porzione di testo selezionata. & \hyperlink{uc24}{UC24}, Capitolato \\ \hline

	% ---- Elementi strutturali ----
	\textbf{R25-F-O} & Obbligatorio & Il sistema deve permettere l'inserimento di intestazioni Markdown (livelli H1--H6) per il testo digitato successivamente. & \hyperlink{uc25}{UC25}, Capitolato, VE-5.1 \\ \hline
	\textbf{R26-F-O} & Obbligatorio & Il sistema deve permettere l'applicazione della formattazione intestazione a una riga di testo già selezionata. & \hyperlink{uc26}{UC26}, Capitolato, VE-5.1 \\ \hline
	\textbf{R27-F-O} & Obbligatorio & Il sistema deve permettere la disattivazione della formattazione intestazione durante la digitazione. & \hyperlink{uc27}{UC27}, Capitolato \\ \hline
	\textbf{R28-F-O} & Obbligatorio & Il sistema deve permettere la rimozione della formattazione intestazione da una riga di testo selezionata, ripristinando il testo normale. & \hyperlink{uc28}{UC28}, Capitolato \\ \hline
	\textbf{R29-F-O} & Obbligatorio & Il sistema deve permettere l'inserimento di un link ipertestuale con testo di visualizzazione e URL. & \hyperlink{uc29}{UC29}, Capitolato, VE-5.1 \\ \hline
	\textbf{R30-F-O} & Obbligatorio & Il sistema deve permettere la modifica di un link ipertestuale esistente. & \hyperlink{uc30}{UC30}, Capitolato \\ \hline
	\textbf{R31-F-O} & Obbligatorio & Il sistema deve permettere la rimozione di un link ipertestuale, mantenendo il testo di visualizzazione come testo normale. & \hyperlink{uc31}{UC31}, Capitolato \\ \hline
	\textbf{R32-F-O} & Obbligatorio & Il sistema deve permettere l'inserimento di un elenco puntato o numerato. & \hyperlink{uc32}{UC32}, \hyperlink{uc33}{UC33}, \hyperlink{uc34}{UC34}, Capitolato, VE-5.1 \\ \hline
	\textbf{R33-F-O} & Obbligatorio & Il sistema deve permettere la modifica di un elenco esistente. & \hyperlink{uc35}{UC35}, \hyperlink{uc36}{UC36}, \hyperlink{uc37}{UC37}, Capitolato \\ \hline
	\textbf{R34-F-O} & Obbligatorio & Il sistema deve permettere la disattivazione dell'elenco durante la digitazione. & \hyperlink{uc38}{UC38}, Capitolato \\ \hline
	\textbf{R35-F-O} & Obbligatorio & Il sistema deve permettere la rimozione della formattazione elenco da una porzione di testo selezionata, ripristinando il testo normale. & \hyperlink{uc39}{UC39}, Capitolato \\ \hline

	% ---- Tabelle ----
	\textbf{R36-F-O} & Obbligatorio & Il sistema deve permettere l'inserimento di una tabella Markdown specificando numero di righe e colonne; i parametri non validi devono essere segnalati con un messaggio di errore. & \hyperlink{uc40}{UC40}, \hyperlink{uc40.1}{UC40.1}, \hyperlink{uc40.2}{UC40.2}, Capitolato, VE-5.1 \\ \hline
	\textbf{R37-F-O} & Obbligatorio & Il sistema deve permettere la modifica del contenuto di una cella di una tabella esistente. & \hyperlink{uc41}{UC41}, Capitolato \\ \hline
	\textbf{R38-F-O} & Obbligatorio & Il sistema deve permettere l'inserimento di una nuova riga in una tabella esistente. & \hyperlink{uc42}{UC42}, Capitolato \\ \hline
	\textbf{R39-F-O} & Obbligatorio & Il sistema deve permettere l'eliminazione di una riga da una tabella esistente. & \hyperlink{uc43}{UC43}, Capitolato \\ \hline
	\textbf{R40-F-O} & Obbligatorio & Il sistema deve permettere l'inserimento di una nuova colonna in una tabella esistente. & \hyperlink{uc44}{UC44}, Capitolato \\ \hline
	\textbf{R41-F-O} & Obbligatorio & Il sistema deve permettere l'eliminazione di una colonna da una tabella esistente. & \hyperlink{uc45}{UC45}, Capitolato \\ \hline
	\textbf{R42-F-O} & Obbligatorio & Il sistema deve permettere l'eliminazione di un'intera tabella dal documento. & \hyperlink{uc46}{UC46}, Capitolato \\ \hline

	% ---- Cronologia modifiche ----
	\textbf{R43-F-O} & Obbligatorio & Il sistema deve permettere l'annullamento dell'ultima modifica apportata al documento (Undo). & \hyperlink{uc47}{UC47}, Analisi del dominio \\ \hline
	\textbf{R44-F-O} & Obbligatorio & Il sistema deve permettere il ripristino dell'ultima modifica annullata (Redo). & \hyperlink{uc48}{UC48}, Analisi del dominio \\ \hline

	% ---- Visualizzazione ----
	\textbf{R45-F-O} & Obbligatorio & Il sistema deve permettere la visualizzazione esclusiva dell'area di scrittura Markdown (Solo Editor). & \hyperlink{uc49}{UC49}, Capitolato \\ \hline
	\textbf{R46-F-O} & Obbligatorio & Il sistema deve permettere la visualizzazione esclusiva dell'anteprima formattata (Solo Render). & \hyperlink{uc50}{UC50}, Capitolato \\ \hline
	\textbf{R47-F-O} & Obbligatorio & Il sistema deve permettere la visualizzazione affiancata di editor e anteprima (Split View). & \hyperlink{uc51}{UC51}, Capitolato \\ \hline

	% ---- Funzionalità AI base ----
	\textbf{R48-F-O} & Obbligatorio & Il sistema deve permettere la generazione di un riassunto del testo selezionato tramite LLM. & \hyperlink{uc52}{UC52}, \hyperlink{uc53}{UC53}, Capitolato \\ \hline
	\textbf{R49-F-O} & Obbligatorio & Il sistema deve notificare l'utente in caso di errore nella comunicazione con i servizi AI durante il riassunto. & \hyperlink{uc60}{UC60}, Analisi del dominio \\ \hline
	\textbf{R50-F-O} & Obbligatorio & Il sistema deve permettere la traduzione del testo selezionato nelle lingue Inglese, Francese e Spagnolo tramite LLM. & \hyperlink{uc61}{UC61}, \hyperlink{uc62}{UC62}, Capitolato, VE-5.2 \\ \hline
	\textbf{R51-F-O} & Obbligatorio & Il sistema deve notificare l'utente in caso di errore nella comunicazione con i servizi AI durante la traduzione. & \hyperlink{uc69}{UC69}, Analisi del dominio \\ \hline
	\textbf{R52-F-D} & Desiderabile & Il sistema dovrebbe supportare la traduzione in Tedesco tramite LLM come estensione facoltativa. & \hyperlink{uc61}{UC61}, \hyperlink{uc62}{UC62}, VE-5.2 \\ \hline
	\textbf{R53-F-O} & Obbligatorio & Il sistema deve permettere la riscrittura del testo selezionato tramite LLM. & \hyperlink{uc70}{UC70}, \hyperlink{uc71}{UC71}, Capitolato \\ \hline
	\textbf{R54-F-O} & Obbligatorio & Il sistema deve notificare l'utente in caso di errore nella comunicazione con i servizi AI durante la riscrittura. & \hyperlink{uc78}{UC78}, Analisi del dominio \\ \hline
	\textbf{R55-F-O} & Obbligatorio & Il sistema deve permettere l'apertura di una finestra dedicata e la generazione di testo a partire da un prompt utente (Distant Writing) tramite LLM. & \hyperlink{uc79}{UC79}, \hyperlink{uc80}{UC80}, \hyperlink{uc81}{UC81}, Capitolato \\ \hline
	\textbf{R56-F-O} & Obbligatorio & Il sistema deve notificare l'utente in caso di errore nella comunicazione con i servizi AI durante la generazione da prompt. & \hyperlink{uc88}{UC88}, Analisi del dominio \\ \hline

	% ---- Sei Cappelli ----
	\textbf{R57-F-O} & Obbligatorio & Il sistema deve permettere l'analisi critica del testo selezionato secondo la prospettiva Cappello Bianco (dati e fatti). & \hyperlink{uc89}{UC89}, \hyperlink{uc90}{UC90}, Capitolato \\ \hline
	\textbf{R58-F-O} & Obbligatorio & Il sistema deve notificare l'utente in caso di errore AI durante l'analisi Cappello Bianco. & \hyperlink{uc97}{UC97}, Analisi del dominio \\ \hline
	\textbf{R59-F-O} & Obbligatorio & Il sistema deve permettere l'analisi critica del testo selezionato secondo la prospettiva Cappello Rosso (emozioni). & \hyperlink{uc98}{UC98}, \hyperlink{uc99}{UC99}, Capitolato \\ \hline
	\textbf{R60-F-O} & Obbligatorio & Il sistema deve notificare l'utente in caso di errore AI durante l'analisi Cappello Rosso. & \hyperlink{uc106}{UC106}, Analisi del dominio \\ \hline
	\textbf{R61-F-O} & Obbligatorio & Il sistema deve permettere l'analisi critica del testo selezionato secondo la prospettiva Cappello Nero (rischi). & \hyperlink{uc107}{UC107}, \hyperlink{uc108}{UC108}, Capitolato \\ \hline
	\textbf{R62-F-O} & Obbligatorio & Il sistema deve notificare l'utente in caso di errore AI durante l'analisi Cappello Nero. & \hyperlink{uc115}{UC115}, Analisi del dominio \\ \hline
	\textbf{R63-F-O} & Obbligatorio & Il sistema deve permettere l'analisi critica del testo selezionato secondo la prospettiva Cappello Giallo (opportunità). & \hyperlink{uc116}{UC116}, \hyperlink{uc117}{UC117}, Capitolato \\ \hline
	\textbf{R64-F-O} & Obbligatorio & Il sistema deve notificare l'utente in caso di errore AI durante l'analisi Cappello Giallo. & \hyperlink{uc124}{UC124}, Analisi del dominio \\ \hline
	\textbf{R65-F-O} & Obbligatorio & Il sistema deve permettere l'analisi critica del testo selezionato secondo la prospettiva Cappello Verde (creatività). & \hyperlink{uc125}{UC125}, \hyperlink{uc126}{UC126}, Capitolato \\ \hline
	\textbf{R66-F-O} & Obbligatorio & Il sistema deve notificare l'utente in caso di errore AI durante l'analisi Cappello Verde. & \hyperlink{uc133}{UC133}, Analisi del dominio \\ \hline
	\textbf{R67-F-O} & Obbligatorio & Il sistema deve permettere l'analisi critica del testo selezionato secondo la prospettiva Cappello Blu (logica/controllo). & \hyperlink{uc134}{UC134}, \hyperlink{uc135}{UC135}, Capitolato \\ \hline
	\textbf{R68-F-O} & Obbligatorio & Il sistema deve notificare l'utente in caso di errore AI durante l'analisi Cappello Blu. & \hyperlink{uc142}{UC142}, Analisi del dominio \\ \hline

	% ---- Gestione proposta AI (requisiti trasversali a tutte le funzionalità AI) ----
	\textbf{R69-F-O} & Obbligatorio & Il sistema deve permettere all'utente di accettare la proposta generata dall'AI, inserendola nel documento. & \hyperlink{uc57}{UC57}, \hyperlink{uc66}{UC66}, \hyperlink{uc75}{UC75}, \hyperlink{uc85}{UC85}, \hyperlink{uc94}{UC94}, \hyperlink{uc103}{UC103}, \hyperlink{uc112}{UC112}, \hyperlink{uc121}{UC121}, \hyperlink{uc130}{UC130}, \hyperlink{uc139}{UC139}, Capitolato, VE-3.2 \\ \hline
	\textbf{R70-F-O} & Obbligatorio & Il sistema deve permettere all'utente di rifiutare la proposta generata dall'AI, mantenendo il documento invariato. & \hyperlink{uc58}{UC58}, \hyperlink{uc67}{UC67}, \hyperlink{uc76}{UC76}, \hyperlink{uc86}{UC86}, \hyperlink{uc95}{UC95}, \hyperlink{uc104}{UC104}, \hyperlink{uc113}{UC113}, \hyperlink{uc122}{UC122}, \hyperlink{uc131}{UC131}, \hyperlink{uc140}{UC140}, Capitolato, VE-3.2 \\ \hline
	\textbf{R71-F-O} & Obbligatorio & Il sistema deve permettere all'utente di richiedere la rigenerazione di una proposta AI non soddisfacente. & \hyperlink{uc59}{UC59}, \hyperlink{uc68}{UC68}, \hyperlink{uc77}{UC77}, \hyperlink{uc87}{UC87}, \hyperlink{uc96}{UC96}, \hyperlink{uc105}{UC105}, \hyperlink{uc114}{UC114}, \hyperlink{uc123}{UC123}, \hyperlink{uc132}{UC132}, \hyperlink{uc141}{UC141}, Analisi del dominio \\ \hline

	% ---- Gestione elaborazione AI (trasversali) ----
	\textbf{R72-F-O} & Obbligatorio & Il sistema deve mostrare all'utente una visualizzazione continua dello stato di avanzamento durante l'elaborazione AI. & \hyperlink{uc54}{UC54}, \hyperlink{uc63}{UC63}, \hyperlink{uc72}{UC72}, \hyperlink{uc82}{UC82}, \hyperlink{uc91}{UC91}, \hyperlink{uc100}{UC100}, \hyperlink{uc109}{UC109}, \hyperlink{uc118}{UC118}, \hyperlink{uc127}{UC127}, \hyperlink{uc136}{UC136}, Analisi del dominio \\ \hline
	\textbf{R73-F-O} & Obbligatorio & Il sistema deve permettere all'utente di interrompere manualmente una richiesta in corso inviata al modello LLM. & \hyperlink{uc55}{UC55}, \hyperlink{uc64}{UC64}, \hyperlink{uc73}{UC73}, \hyperlink{uc83}{UC83}, \hyperlink{uc92}{UC92}, \hyperlink{uc101}{UC101}, \hyperlink{uc110}{UC110}, \hyperlink{uc119}{UC119}, \hyperlink{uc128}{UC128}, \hyperlink{uc137}{UC137}, Verbale Esterno \\ \hline
	\textbf{R74-F-O} & Obbligatorio & Il sistema deve visualizzare la proposta generata dall'AI all'interno di un modale dedicato, separato dal testo della nota. & \hyperlink{uc56}{UC56}, \hyperlink{uc65}{UC65}, \hyperlink{uc74}{UC74}, \hyperlink{uc84}{UC84}, \hyperlink{uc93}{UC93}, \hyperlink{uc102}{UC102}, \hyperlink{uc111}{UC111}, \hyperlink{uc120}{UC120}, \hyperlink{uc129}{UC129}, \hyperlink{uc138}{UC138}, Capitolato \\ \hline

	% ---- I/O e persistenza locale ----
	\textbf{R75-F-O} & Obbligatorio & Il sistema deve permettere il salvataggio della nota corrente in un file \texttt{.md} nel file system locale; deve notificare in caso di errore. & \hyperlink{uc145}{UC145}, \hyperlink{uc145.1}{UC145.1}, Capitolato \\ \hline
	\textbf{R76-F-O} & Obbligatorio & Il sistema deve permettere il caricamento di un file \texttt{.md} esistente dal file system locale nell'editor; deve notificare in caso di file non valido o corrotto. & \hyperlink{uc146}{UC146}, \hyperlink{uc146.1}{UC146.1}, Capitolato \\ \hline
	\textbf{R77-F-O} & Obbligatorio & Il sistema deve permettere l'esportazione della nota nei formati PDF, HTML e JSON; deve notificare in caso di errore durante la conversione. & \hyperlink{uc147}{UC147}, \hyperlink{uc147.1}{UC147.1}, \hyperlink{uc148}{UC148}, \hyperlink{uc148.1}{UC148.1}, \hyperlink{uc149}{UC149}, \hyperlink{uc149.1}{UC149.1}, Capitolato \\ \hline

	% ---- Opzionali: DB remoto ----
	\textbf{R78-F-Op} & Opzionale & Il sistema deve permettere la creazione di una nuova nota vuota associando il titolo inserito dall'utente. & \hyperlink{uc143}{UC143}, Capitolato \\ \hline
	\textbf{R79-F-Op} & Opzionale & Il sistema deve permettere l'eliminazione di una nota previa conferma esplicita da parte dell'utente. & \hyperlink{uc144}{UC144}, \hyperlink{uc144.1}{UC144.1}, \hyperlink{uc144.2}{UC144.2}, Capitolato \\ \hline
	\textbf{R80-F-Op} & Opzionale & Il sistema deve notificare l'utente in caso di errore tecnico o del file system durante l'eliminazione di una nota. & \hyperlink{uc144.3}{UC144.3}, Analisi del dominio \\ \hline
	\textbf{R81-F-Op} & Opzionale & Il sistema deve permettere la registrazione di un nuovo utente tramite email e password, con validazione dei dati inseriti. & \hyperlink{uc150}{UC150}, \hyperlink{uc150.1}{UC150.1}, \hyperlink{uc150.2}{UC150.2}, Capitolato \\ \hline
	\textbf{R82-F-Op} & Opzionale & Il sistema deve permettere l'autenticazione dell'utente tramite credenziali; deve notificare in caso di credenziali errate. & \hyperlink{uc151}{UC151}, \hyperlink{uc151.1}{UC151.1}, Capitolato \\ \hline
	\textbf{R83-F-Op} & Opzionale & Il sistema deve permettere il logout dell'utente autenticato. & \hyperlink{uc152}{UC152}, Capitolato \\ \hline
	\textbf{R84-F-Op} & Opzionale & Il sistema deve permettere la visualizzazione della lista delle note salvate nel servizio remoto con nome e data di ultima modifica; deve notificare in caso di errore. & \hyperlink{uc153}{UC153}, \hyperlink{uc153.1}{UC153.1}, \hyperlink{uc153.1.1}{UC153.1.1}, \hyperlink{uc153.1.2}{UC153.1.2}, \hyperlink{uc153.1.3}{UC153.1.3}, \hyperlink{uc153.2}{UC153.2}, Capitolato \\ \hline
	\textbf{R85-F-Op} & Opzionale & Il sistema deve permettere il salvataggio della nota corrente nello spazio remoto associato al profilo utente; deve notificare in caso di errore di sincronizzazione. & \hyperlink{uc154}{UC154}, \hyperlink{uc154.1}{UC154.1}, Capitolato \\ \hline
	\textbf{R86-F-Op} & Opzionale & Il sistema deve permettere l'apertura di una nota dalla lista delle note salvate nel servizio remoto; deve notificare in caso di errore nel caricamento. & \hyperlink{uc155}{UC155}, \hyperlink{uc155.1}{UC155.1}, Capitolato \\ \hline

\end{longtable}

% -----------------------------------------------------------------------
\subsection{Requisiti di qualità}
\begin{longtable}{|p{2.5cm}|p{2.2cm}|p{6.8cm}|p{2.7cm}|}
	\caption{Requisiti di Qualità} \\
	\hline
	\rowcolor{primarycolor!10}
	\textbf{Codice} & \textbf{Importanza} & \textbf{Descrizione} & \textbf{UC, Fonti} \\
	\hline
	\endfirsthead
	\hline
	\rowcolor{primarycolor!10}
	\textbf{Codice} & \textbf{Importanza} & \textbf{Descrizione} & \textbf{UC, Fonti} \\
	\hline
	\endhead

	\textbf{R1-Q-O} & Obbligatorio & Il sistema deve essere corredato da una suite di test automatici che garantisca una line coverage minima dell'80\% sui moduli che implementano i casi d'uso principali, verificabile tramite il report di copertura generato dal framework di test adottato. & Capitolato \\ \hline
	\textbf{R2-Q-O} & Obbligatorio & È necessario produrre il documento di Specifica Tecnica, contenente i diagrammi e la descrizione dell'architettura del software. & Capitolato \\ \hline
	\textbf{R3-Q-O} & Obbligatorio & È necessario produrre una documentazione dettagliata riguardante i test previsti ed eseguiti (Piano di Qualifica). & Capitolato \\ \hline
	\textbf{R4-Q-O} & Obbligatorio & È necessario produrre il documento "Manuale Utente", con indicazioni chiare riguardanti l'utilizzo del software lato utente. & Capitolato \\ \hline
	\textbf{R5-Q-O} & Obbligatorio & Il progetto deve essere sviluppato utilizzando Git\textsubscript{G} come strumento di versionamento del codice, con tracciamento delle attività tramite issue. & Interno \\ \hline
	\textbf{R6-Q-O} & Obbligatorio & È necessario eseguire test di integrazione che verifichino la corretta comunicazione tra l'applicazione e i servizi LLM esterni, assicurando che le richieste e le risposte vengano scambiate correttamente. & Analisi del dominio \\ \hline
	\textbf{R7-Q-Op} & Opzionale & È necessario eseguire test di integrazione che verifichino la corretta sincronizzazione dei dati tra il client e il servizio di persistenza remota, qualora il modulo opzionale sia implementato. & Analisi del dominio \\ \hline
	\textbf{R8-Q-Op} & Opzionale & È necessario eseguire test che verifichino la corretta applicazione dei meccanismi di autenticazione e autorizzazione nel modulo opzionale di gestione utente. & Analisi del dominio \\ \hline
	\textbf{R9-Q-D} & Desiderabile & Il progetto dovrebbe adottare una licenza open source permissiva (ad esempio MIT o Apache License 2.0), in linea con le indicazioni dell'azienda proponente sulla natura open source del software rilasciato. & Capitolato \\ \hline

\end{longtable}

% -----------------------------------------------------------------------
\subsection{Requisiti di vincolo}
\begin{longtable}{|p{2.5cm}|p{2.2cm}|p{6.8cm}|p{2.7cm}|}
	\caption{Requisiti di Vincolo} \\
	\hline
	\rowcolor{primarycolor!10}
	\textbf{Codice} & \textbf{Importanza} & \textbf{Descrizione} & \textbf{UC, Fonti} \\
	\hline
	\endfirsthead
	\hline
	\rowcolor{primarycolor!10}
	\textbf{Codice} & \textbf{Importanza} & \textbf{Descrizione} & \textbf{UC, Fonti} \\
	\hline
	\endhead

	\textbf{R1-V-O} & Obbligatorio & L'applicazione deve essere eseguibile direttamente all'interno di un browser web standard, basata su tecnologie HTML5, CSS3 e JavaScript/TypeScript. & Capitolato \\ \hline
	\textbf{R2-V-O} & Obbligatorio & L'interazione con i modelli LLM deve avvenire tramite chiamate API aderenti allo standard OpenAI. & Capitolato, VE-1.1 \\ \hline
	\textbf{R3-V-O} & Obbligatorio & I modelli LLM di riferimento sono Gemma, ChatGPT-OSS e Qwen, messi a disposizione dall'azienda proponente; l'infrastruttura attuale prevede l'utilizzo iniziale del modello \texttt{gemma4:12b}, con possibile ampliamento dei permessi di accesso al modello \texttt{Qwen3.6-35B-A3B}, ritenuto più stabile e performante, qualora necessario. & Capitolato, VE-1.1, VE-7.2 \\ \hline
	\textbf{R4-V-Op} & Opzionale & Il sistema utilizza un database PostgreSQL\textsubscript{G} configurato a livello applicativo per la memorizzazione remota delle note. & VE-1.4, VE-7.3 \\ \hline
	\textbf{R5-V-O} & Obbligatorio & Le componenti web del sistema devono essere compatibili con le versioni di Chrome 151+, Edge 150+, Firefox 153+ e Safari 26+. & VE-7.3 \\ \hline
	\textbf{R6-V-O} & Obbligatorio & Ogni richiesta di elaborazione AI deve essere indirizzata a un singolo modello LLM per l'intera durata della richiesta, poiché l'infrastruttura prevede il ricaricamento del modello ad ogni chiamata: concatenare chiamate a modelli differenti comporterebbe una perdita di efficienza. & VE-7.2 \\ \hline
	\textbf{R7-V-O} & Obbligatorio & Il codice sorgente e la documentazione del progetto devono essere pubblicati su un repository pubblico. & Capitolato, VE-3.3, VE-5.3 \\ \hline

\end{longtable}

% -----------------------------------------------------------------------
\subsection{Requisiti prestazionali}
\begin{longtable}{|p{2.5cm}|p{2.2cm}|p{6.8cm}|p{2.7cm}|}
	\caption{Requisiti Prestazionali} \\
	\hline
	\rowcolor{primarycolor!10}
	\textbf{Codice} & \textbf{Importanza} & \textbf{Descrizione} & \textbf{UC, Fonti} \\
	\hline
	\endfirsthead
	\hline
	\rowcolor{primarycolor!10}
	\textbf{Codice} & \textbf{Importanza} & \textbf{Descrizione} & \textbf{UC, Fonti} \\
	\hline
	\endhead

	\textbf{R1-P-O} & Obbligatorio & In condizioni di funzionamento standard (dispositivo conforme ai requisiti hardware minimi e documento di dimensioni ordinarie), il rendering dell'anteprima Markdown deve aggiornarsi in tempo reale durante la digitazione, entro 1 secondo dalla modifica del contenuto dell'editor. & Capitolato, Analisi del dominio \\ \hline
	\textbf{R2-P-O} & Obbligatorio & A fronte di una connessione di rete stabile verso le API, il sistema deve garantire una visualizzazione continua e reattiva dello stato di caricamento per l'intera durata dell'elaborazione AI, indipendentemente dai tempi di risposta del modello, mantenendo l'indicatore di caricamento visibile per un tempo minimo di circa 2 secondi. La durata dell'elaborazione è determinata dalla risposta del modello e può essere interrotta in qualsiasi momento dall'utente tramite il meccanismo di interruzione manuale (R73-F-O). & Analisi del dominio, VE-6.1, VE-7.3 \\ \hline

\end{longtable}